\section{Conclusion}
\label{sec:conclusion}

The event-level compressor case studies gave a direct negative answer. MetroPT contained one independent labelled test failure; the registered diagnostic detected it, but only with 2527 predicted alarm episodes, event precision \(0.00040\), and a false-alarm burden of 101.17 alarm episodes per operating day. MetroPT2 also contained one independent labelled test failure; it was detected with 849 predicted alarm episodes, event precision \(0.00118\), and 51.29 false-alarm episodes per operating day. These detections are therefore not operationally persuasive: the event was recovered, but the alarm burden was far beyond what would constitute a useful maintenance signal.

The statistical interpretation is that clustered extremes are dependence diagnostics, not maintenance decisions. Threshold and run-length choices changed the estimated clustering behavior, which means the extremal-index estimate is conditional on the observable, regime definition, threshold, and declustering rule. Low estimated extremal index can indicate recurrence, persistence, smoothing, missingness, periodic control behavior, or regime mixture; it is not direct evidence of degradation unless the target region is known to be failure-proximal and the episode-level alarm policy remains useful after conversion.

The four research questions are answered narrowly. Regime conditioning and threshold grids exposed sensitivity, but they did not establish a stable diagnostic suitable for general deployment. Declustering and alarm merging did not reduce duplicate alarms while preserving event recall in the registered failed-detection grid; all tested merge-gap and post-onset tolerance combinations had zero recall, precision, and \(F_1\). The robust extreme-score diagnostic added limited diagnostic value by revealing alarm-conversion failure and estimand mismatch, but it did not add operational early-warning value. The dominant failure mechanisms are sparse independent failures and alarm burden for MetroPT and MetroPT2, vehicle-level estimand mismatch for SCANIA Component X, cycle-state mismatch for Hydraulic Systems, and yield-target mismatch for SECOM.

Recurrence-based EVT can support maintenance decisions only after stronger conditions are met: an identifiable failure-proximal target region, enough independent failures or entities for uncertainty estimation, leakage-free validation, stable regimes and thresholds, fair comparison against simple baselines, acceptable event recall, acceptable false-alarm burden, useful warning lead time, and prospective or external validation. The present study should therefore be read as evidence about the limits of converting clustered extremes into maintenance alarms, not as evidence of predictive superiority.
